% =============================================================================
\section{Project Overview}
% =============================================================================

\subsection{Purpose}

The \texttt{star-gateway} is the central backend service running on the Raspberry Pi 5. It provides:

\begin{itemize}
    \item \textbf{GraphQL API} for telemetry, settings, and navigation queries
    \item \textbf{gRPC service} for low-latency motor commands
    \item \textbf{Hardware adapters} for ESP32 communication (SPI + TCP)
    \item \textbf{ROS 2 bridge} for SLAM and navigation integration
    \item \textbf{Prometheus metrics} for monitoring
\end{itemize}

\subsection{Technology Stack}

\begin{table}[H]
\centering
\begin{tabular}{|l|l|l|}
\hline
\textbf{Component} & \textbf{Technology} & \textbf{Version} \\
\hline
Language & Kotlin & 1.9.22 \\
Framework & Spring Boot & 3.2.x \\
GraphQL & Spring GraphQL & 1.2.x \\
gRPC & grpc-spring-boot-starter & 3.1.x \\
Hardware Access & Pi4J & 2.6.x \\
Build Tool & Gradle (Kotlin DSL) & 8.x \\
JVM & OpenJDK & 17 \\
\hline
\end{tabular}
\caption{Technology Stack}
\end{table}

\subsection{Architecture Overview}

\begin{figure}[H]
\centering
\begin{tikzpicture}[
    node distance=1.5cm,
    box/.style={draw, rectangle, rounded corners, minimum width=3cm, minimum height=1cm, align=center, font=\small},
    arrow/.style={->, thick}
]
    % External clients
    \node[box, fill=green!20] (pwa) {star-ui (PWA)};
    \node[box, fill=blue!20, right=3cm of pwa] (grpcweb) {gRPC-Web\\(Envoy)};

    % Gateway
    \node[box, fill=orange!20, below=2cm of pwa, minimum width=8cm] (gateway) {star-gateway (Spring Boot)};

    % Internal components
    \node[box, fill=yellow!20, below left=2cm and 1cm of gateway] (esp32) {ESP32\\(SPI/TCP)};
    \node[box, fill=purple!20, below=2cm of gateway] (ros2) {ROS 2\\(DDS)};
    \node[box, fill=red!20, below right=2cm and 1cm of gateway] (lidar) {LiDAR\\(rplidar\_ros2)};

    % Connections
    \draw[arrow, <->] (pwa) -- node[left] {GraphQL} (gateway);
    \draw[arrow, <->] (grpcweb) -- node[right] {gRPC} (gateway);
    \draw[arrow, <->] (gateway) -- node[left] {SPI/TCP} (esp32);
    \draw[arrow, <->] (gateway) -- node[right] {DDS} (ros2);
    \draw[arrow] (lidar) -- (ros2);
\end{tikzpicture}
\caption{Gateway Architecture}
\end{figure}

% =============================================================================
\section{Project Structure}
% =============================================================================

\begin{lstlisting}[language=bash]
star-gateway/
|-- build.gradle.kts              # Build configuration
|-- settings.gradle.kts           # Project settings
|-- src/
|   +-- main/
|       |-- kotlin/com/star/gateway/
|       |   |-- StarGatewayApplication.kt
|       |   |-- config/
|       |   |   |-- GraphQLConfig.kt
|       |   |   |-- GrpcConfig.kt
|       |   |   +-- CorsConfig.kt
|       |   |-- graphql/
|       |   |   |-- TelemetryController.kt
|       |   |   |-- SettingsController.kt
|       |   |   +-- NavigationController.kt
|       |   |-- grpc/
|       |   |   +-- MotorControlGrpcService.kt
|       |   |-- service/
|       |   |   |-- MotorControlService.kt
|       |   |   |-- TelemetryService.kt
|       |   |   |-- NavigationService.kt
|       |   |   +-- SensorFusionService.kt
|       |   |-- hardware/
|       |   |   |-- Esp32SpiAdapter.kt
|       |   |   |-- Esp32TcpAdapter.kt
|       |   |   |-- RplidarAdapter.kt
|       |   |   +-- Ros2Bridge.kt
|       |   +-- model/
|       |       |-- TelemetryData.kt
|       |       |-- MotorCommand.kt
|       |       +-- NavigationState.kt
|       |-- proto/
|       |   +-- motor_control.proto
|       +-- resources/
|           |-- application.yml
|           +-- graphql/
|               +-- schema.graphqls
+-- deployment/
    |-- star-gateway.service      # systemd unit
    +-- envoy.yaml                # gRPC-Web proxy config
\end{lstlisting}

% =============================================================================
\section{Implementation Tasks}
% =============================================================================

\subsection{Phase 1: Project Setup}

\subsubsection{Task 1.1: Gradle Build Configuration}

Create \texttt{build.gradle.kts}:

\begin{lstlisting}[language=Kotlin, caption=build.gradle.kts]
plugins {
    kotlin("jvm") version "1.9.22"
    kotlin("plugin.spring") version "1.9.22"
    id("org.springframework.boot") version "3.2.4"
    id("io.spring.dependency-management") version "1.1.4"
    id("com.google.protobuf") version "0.9.4"
}

java {
    sourceCompatibility = JavaVersion.VERSION_17
}

repositories {
    mavenCentral()
}

dependencies {
    // Spring Boot Core
    implementation("org.springframework.boot:spring-boot-starter-webflux")
    implementation("org.springframework.boot:spring-boot-starter-actuator")

    // GraphQL
    implementation("org.springframework.boot:spring-boot-starter-graphql")

    // gRPC
    implementation("net.devh:grpc-spring-boot-starter:3.1.0.RELEASE")
    implementation("io.grpc:grpc-kotlin-stub:1.4.1")
    implementation("io.grpc:grpc-protobuf:1.60.1")
    implementation("com.google.protobuf:protobuf-kotlin:3.25.2")

    // Kotlin Coroutines
    implementation("org.jetbrains.kotlinx:kotlinx-coroutines-core:1.8.0")
    implementation("org.jetbrains.kotlinx:kotlinx-coroutines-reactor:1.8.0")

    // Pi4J for Hardware Access (RPi5 support)
    implementation("com.pi4j:pi4j-core:2.6.1")
    implementation("com.pi4j:pi4j-plugin-raspberrypi:2.6.1")
    implementation("com.pi4j:pi4j-plugin-gpiod:2.6.1")
    implementation("com.pi4j:pi4j-plugin-linuxfs:2.6.1")

    // Metrics
    implementation("io.micrometer:micrometer-registry-prometheus")

    // JSON
    implementation("com.fasterxml.jackson.module:jackson-module-kotlin")

    // Kotlin
    implementation("org.jetbrains.kotlin:kotlin-reflect")

    // Testing
    testImplementation("org.springframework.boot:spring-boot-starter-test")
}

protobuf {
    protoc {
        artifact = "com.google.protobuf:protoc:3.25.2"
    }
    plugins {
        id("grpc") {
            artifact = "io.grpc:protoc-gen-grpc-java:1.60.1"
        }
        id("grpckt") {
            artifact = "io.grpc:protoc-gen-grpc-kotlin:1.4.1:jdk8@jar"
        }
    }
    generateProtoTasks {
        all().forEach {
            it.plugins {
                id("grpc")
                id("grpckt")
            }
            it.builtins {
                id("kotlin")
            }
        }
    }
}
\end{lstlisting}

\subsubsection{Task 1.2: Application Configuration}

Create \texttt{src/main/resources/application.yml}:

\begin{lstlisting}[language=yaml, caption=application.yml]
spring:
  application:
    name: star-gateway
  graphql:
    graphiql:
      enabled: true
      path: /graphiql
    websocket:
      path: /graphql-ws
    schema:
      locations: classpath:graphql/

server:
  port: 8080

grpc:
  server:
    port: 9090

star:
  esp32:
    spi:
      enabled: true
      bus: 0
      chip-select: 0
      speed-hz: 10000000
    tcp:
      enabled: true
      host: 192.168.4.2
      port: 5000
      timeout-ms: 100
  lidar:
    enabled: true
    topic: /scan

management:
  endpoints:
    web:
      exposure:
        include: health,prometheus,info,metrics
  endpoint:
    health:
      show-details: always

logging:
  level:
    com.star.gateway: DEBUG
\end{lstlisting}

\subsection{Phase 2: GraphQL Implementation}

\subsubsection{Task 2.1: GraphQL Schema}

Create \texttt{src/main/resources/graphql/schema.graphqls}:

\begin{lstlisting}[language=GraphQL, caption=schema.graphqls]
type Query {
    telemetry: Telemetry!
    systemStatus: SystemStatus!
    settings: Settings!
    navigationState: NavigationState!
    telemetryHistory(limit: Int = 100): [TelemetrySnapshot!]!
}

type Mutation {
    setMode(mode: RobotMode!): Boolean!
    emergencyStop: Boolean!
    setSettings(input: SettingsInput!): Settings!
    addWaypoint(x: Float!, y: Float!): Waypoint!
    clearWaypoints: Boolean!
    startNavigation: Boolean!
    stopNavigation: Boolean!
}

type Subscription {
    telemetryStream: Telemetry!
    lidarScan: LidarScan!
    eventLog: Event!
}

type Telemetry {
    imu: ImuData!
    gps: GpsData!
    battery: Float!
    wifiSignal: Int!
    cpuUsage: Float!
    temperature: Float!
    motorLoad: Float!
    timestamp: String!
}

type ImuData {
    pitch: Float!
    roll: Float!
    yaw: Float!
    velocity: Float!
    accelerationX: Float!
    accelerationY: Float!
    accelerationZ: Float!
}

type SystemStatus {
    connectionStatus: ConnectionStatus!
    mode: RobotMode!
    esp32Connected: Boolean!
    lidarConnected: Boolean!
    rosConnected: Boolean!
}

type NavigationState {
    currentPosition: Position!
    path: [Position!]!
    waypoints: [Waypoint!]!
    isNavigating: Boolean!
}

enum RobotMode {
    IDLE
    MANUAL
    AUTONOMOUS
    MAPPING
    EMERGENCY_STOP
}

enum ConnectionStatus {
    CONNECTED
    DISCONNECTED
    CONNECTING
    ERROR
}

input SettingsInput {
    maxSpeed: Float
    safetyLockEnabled: Boolean
    autonomousEnabled: Boolean
}
\end{lstlisting}

\subsubsection{Task 2.2: GraphQL Controller}

Create \texttt{src/main/kotlin/com/star/gateway/graphql/TelemetryController.kt}:

\begin{lstlisting}[language=Kotlin, caption=TelemetryController.kt]
package com.star.gateway.graphql

import com.star.gateway.model.*
import com.star.gateway.service.TelemetryService
import com.star.gateway.service.MotorControlService
import kotlinx.coroutines.flow.Flow
import org.springframework.graphql.data.method.annotation.*
import org.springframework.stereotype.Controller
import reactor.core.publisher.Flux

@Controller
class TelemetryController(
    private val telemetryService: TelemetryService,
    private val motorControlService: MotorControlService
) {
    @QueryMapping
    suspend fun telemetry(): Telemetry {
        return telemetryService.getCurrentTelemetry()
    }

    @QueryMapping
    suspend fun systemStatus(): SystemStatus {
        return telemetryService.getSystemStatus()
    }

    @MutationMapping
    suspend fun emergencyStop(): Boolean {
        val result = motorControlService.emergencyStop()
        return result.success
    }

    @MutationMapping
    suspend fun setMode(@Argument mode: RobotMode): Boolean {
        return telemetryService.setMode(mode)
    }

    @SubscriptionMapping
    fun telemetryStream(): Flux<Telemetry> {
        return telemetryService.telemetryFlow().asFlux()
    }
}
\end{lstlisting}

\subsection{Phase 3: gRPC Implementation}

\subsubsection{Task 3.1: Protobuf Definition}

Create \texttt{src/main/proto/motor\_control.proto}:

\begin{lstlisting}[language=Protobuf, caption=motor\_control.proto]
syntax = "proto3";

package star.motor;

option java_package = "com.star.gateway.proto";
option java_multiple_files = true;

service MotorControl {
    rpc SetVelocity(VelocityCommand) returns (CommandResponse);
    rpc EmergencyStop(EmptyRequest) returns (CommandResponse);
    rpc StreamEncoderData(EmptyRequest) returns (stream EncoderData);
}

message VelocityCommand {
    float left_velocity = 1;
    float right_velocity = 2;
    uint64 timestamp = 3;
}

message CommandResponse {
    bool success = 1;
    string message = 2;
    uint64 latency_us = 3;
}

message EncoderData {
    int32 motor_id = 1;
    int64 ticks = 2;
    float velocity = 3;
    uint64 timestamp = 4;
}

message EmptyRequest {}
\end{lstlisting}

\subsubsection{Task 3.2: gRPC Service}

Create \texttt{src/main/kotlin/com/star/gateway/grpc/MotorControlGrpcService.kt}:

\begin{lstlisting}[language=Kotlin, caption=MotorControlGrpcService.kt]
package com.star.gateway.grpc

import com.star.gateway.proto.*
import com.star.gateway.service.MotorControlService
import kotlinx.coroutines.flow.Flow
import kotlinx.coroutines.flow.flow
import net.devh.boot.grpc.server.service.GrpcService
import org.slf4j.LoggerFactory

@GrpcService
class MotorControlGrpcService(
    private val motorControlService: MotorControlService
) : MotorControlGrpcKt.MotorControlCoroutineImplBase() {

    private val logger = LoggerFactory.getLogger(javaClass)

    override suspend fun setVelocity(request: VelocityCommand): CommandResponse {
        val startTime = System.nanoTime()

        logger.debug("gRPC SetVelocity: left=${request.leftVelocity}, right=${request.rightVelocity}")

        val result = motorControlService.setVelocity(
            left = request.leftVelocity,
            right = request.rightVelocity
        )

        val latencyUs = (System.nanoTime() - startTime) / 1000

        return CommandResponse.newBuilder()
            .setSuccess(result.success)
            .setMessage(result.message)
            .setLatencyUs(latencyUs)
            .build()
    }

    override suspend fun emergencyStop(request: EmptyRequest): CommandResponse {
        logger.warn("gRPC EmergencyStop received")
        val result = motorControlService.emergencyStop()

        return CommandResponse.newBuilder()
            .setSuccess(result.success)
            .setMessage(result.message)
            .build()
    }

    override fun streamEncoderData(request: EmptyRequest): Flow<EncoderData> = flow {
        while (true) {
            val encoders = motorControlService.getEncoderData()
            encoders.forEach { (motorId, data) ->
                emit(
                    EncoderData.newBuilder()
                        .setMotorId(motorId)
                        .setTicks(data.ticks)
                        .setVelocity(data.velocity)
                        .setTimestamp(System.currentTimeMillis())
                        .build()
                )
            }
            kotlinx.coroutines.delay(10) // 100Hz
        }
    }
}
\end{lstlisting}

\subsection{Phase 4: Hardware Adapters}

\subsubsection{Task 4.1: ESP32 SPI Adapter}

Create \texttt{src/main/kotlin/com/star/gateway/hardware/Esp32SpiAdapter.kt}:

\begin{lstlisting}[language=Kotlin, caption=Esp32SpiAdapter.kt]
package com.star.gateway.hardware

import com.pi4j.Pi4J
import com.pi4j.context.Context
import com.pi4j.io.spi.Spi
import com.pi4j.io.spi.SpiMode
import com.star.gateway.model.CommandResult
import org.springframework.beans.factory.annotation.Value
import org.springframework.stereotype.Component
import java.nio.ByteBuffer
import java.nio.ByteOrder
import javax.annotation.PostConstruct
import javax.annotation.PreDestroy

@Component
class Esp32SpiAdapter(
    @Value("\${star.esp32.spi.enabled:true}")
    private val enabled: Boolean,
    @Value("\${star.esp32.spi.bus:0}")
    private val spiBus: Int,
    @Value("\${star.esp32.spi.chip-select:0}")
    private val chipSelect: Int,
    @Value("\${star.esp32.spi.speed-hz:10000000}")
    private val speedHz: Int
) {
    companion object {
        const val CMD_SET_VELOCITY: Byte = 0x01
        const val CMD_EMERGENCY_STOP: Byte = 0x02
        const val CMD_GET_ENCODERS: Byte = 0x03
        const val RESPONSE_OK: Byte = 0x00
    }

    private var pi4j: Context? = null
    private var spi: Spi? = null
    private var connected = false

    @PostConstruct
    fun initialize() {
        if (!enabled) return

        try {
            pi4j = Pi4J.newAutoContext()
            val spiConfig = Spi.newConfigBuilder(pi4j)
                .id("esp32-spi")
                .bus(spiBus)
                .chipSelect(chipSelect)
                .baud(speedHz)
                .mode(SpiMode.MODE_0)
                .build()

            spi = pi4j?.create(spiConfig)
            connected = testConnection()
        } catch (e: Exception) {
            connected = false
        }
    }

    @PreDestroy
    fun shutdown() {
        spi?.close()
        pi4j?.shutdown()
    }

    suspend fun sendVelocityCommand(left: Float, right: Float): CommandResult {
        if (!enabled || spi == null) {
            throw IllegalStateException("SPI not initialized")
        }

        val buffer = ByteBuffer.allocate(9)
            .order(ByteOrder.LITTLE_ENDIAN)
            .put(CMD_SET_VELOCITY)
            .putFloat(left)
            .putFloat(right)

        val txData = buffer.array()
        val rxData = ByteArray(2)

        synchronized(this) {
            spi!!.transfer(txData, rxData)
        }

        return if (rxData[0] == RESPONSE_OK) {
            CommandResult.success("Velocity set: left=$left, right=$right")
        } else {
            CommandResult.failure("ESP32 error: ${rxData[0]}")
        }
    }

    suspend fun emergencyStop(): CommandResult {
        val txData = byteArrayOf(CMD_EMERGENCY_STOP)
        val rxData = ByteArray(2)

        synchronized(this) {
            spi?.transfer(txData, rxData)
        }

        return if (rxData[0] == RESPONSE_OK) {
            CommandResult.success("Emergency stop executed")
        } else {
            CommandResult.failure("Emergency stop failed")
        }
    }

    private fun testConnection(): Boolean {
        // Ping command implementation
        return true
    }

    fun isConnected(): Boolean = connected && enabled
}
\end{lstlisting}

\subsubsection{Task 4.2: ESP32 TCP Adapter}

Create \texttt{src/main/kotlin/com/star/gateway/hardware/Esp32TcpAdapter.kt} as fallback communication channel using JSON over TCP socket.

\subsubsection{Task 4.3: ROS 2 Bridge}

Create \texttt{src/main/kotlin/com/star/gateway/hardware/Ros2Bridge.kt} to interface with ROS 2 via DDS or rosbridge.

\subsection{Phase 5: Deployment}

\subsubsection{Task 5.1: systemd Service}

Create \texttt{deployment/star-gateway.service}:

\begin{lstlisting}[language=ini, caption=star-gateway.service]
[Unit]
Description=STAR Robot Gateway
After=network-online.target
Wants=network-online.target

[Service]
Type=simple
User=root
WorkingDirectory=/opt/star-gateway
ExecStart=/usr/bin/java -Xmx512m -jar star-gateway.jar
Environment="JAVA_HOME=/usr/lib/jvm/java-17-openjdk"
Restart=always
RestartSec=10

[Install]
WantedBy=multi-user.target
\end{lstlisting}

\subsubsection{Task 5.2: Envoy Proxy (gRPC-Web)}

Create \texttt{deployment/envoy.yaml} for browser gRPC support:

\begin{lstlisting}[language=yaml, caption=envoy.yaml]
static_resources:
  listeners:
  - name: listener_0
    address:
      socket_address: { address: 0.0.0.0, port_value: 8090 }
    filter_chains:
    - filters:
      - name: envoy.filters.network.http_connection_manager
        typed_config:
          "@type": type.googleapis.com/envoy.extensions.filters.network.http_connection_manager.v3.HttpConnectionManager
          codec_type: auto
          stat_prefix: ingress_http
          route_config:
            name: local_route
            virtual_hosts:
            - name: local_service
              domains: ["*"]
              routes:
              - match: { prefix: "/" }
                route:
                  cluster: grpc_service
                  timeout: 0s
                  max_stream_duration:
                    grpc_timeout_header_max: 0s
              cors:
                allow_origin_string_match:
                - prefix: "*"
                allow_methods: GET, PUT, DELETE, POST, OPTIONS
                allow_headers: keep-alive,user-agent,cache-control,content-type,content-transfer-encoding,x-accept-content-transfer-encoding,x-accept-response-streaming,x-user-agent,x-grpc-web,grpc-timeout
                expose_headers: grpc-status,grpc-message
          http_filters:
          - name: envoy.filters.http.grpc_web
          - name: envoy.filters.http.cors
          - name: envoy.filters.http.router

  clusters:
  - name: grpc_service
    connect_timeout: 0.25s
    type: logical_dns
    http2_protocol_options: {}
    lb_policy: round_robin
    load_assignment:
      cluster_name: cluster_0
      endpoints:
      - lb_endpoints:
        - endpoint:
            address:
              socket_address:
                address: localhost
                port_value: 9090
\end{lstlisting}

% =============================================================================
\section{Build and Run}
% =============================================================================

\subsection{Local Development}

\begin{lstlisting}[language=bash]
cd star-gateway

# Build
./gradlew build

# Run locally (without Pi4J hardware)
./gradlew bootRun --args='--star.esp32.spi.enabled=false'

# Build JAR for deployment
./gradlew bootJar

# Output: build/libs/star-gateway-*.jar
\end{lstlisting}

\subsection{Deploy to Pi5}

\begin{lstlisting}[language=bash]
# Copy JAR to Pi5
scp build/libs/star-gateway-*.jar root@star-robot.local:/opt/star-gateway/star-gateway.jar

# SSH and restart service
ssh root@star-robot.local
systemctl restart star-gateway
journalctl -u star-gateway -f
\end{lstlisting}

\subsection{API Endpoints}

\begin{table}[H]
\centering
\begin{tabular}{|l|l|l|}
\hline
\textbf{Endpoint} & \textbf{Protocol} & \textbf{Purpose} \\
\hline
\texttt{http://star-robot.local:8080/graphql} & GraphQL & Queries, Mutations \\
\texttt{ws://star-robot.local:8080/graphql-ws} & WebSocket & Subscriptions \\
\texttt{http://star-robot.local:8080/graphiql} & HTTP & GraphQL IDE \\
\texttt{star-robot.local:9090} & gRPC & Motor commands \\
\texttt{star-robot.local:8090} & gRPC-Web & Browser gRPC \\
\texttt{http://star-robot.local:8080/actuator/*} & HTTP & Health, Metrics \\
\hline
\end{tabular}
\caption{API Endpoints}
\end{table}

% =============================================================================
\section{Testing}
% =============================================================================

\subsection{GraphQL Testing}

Use GraphiQL at \texttt{http://star-robot.local:8080/graphiql}:

\begin{lstlisting}[language=GraphQL, caption=Sample GraphQL Queries]
# Get current telemetry
query {
  telemetry {
    battery
    temperature
    imu { pitch roll yaw }
  }
}

# Emergency stop
mutation {
  emergencyStop
}

# Subscribe to telemetry stream
subscription {
  telemetryStream {
    battery
    timestamp
  }
}
\end{lstlisting}

\subsection{gRPC Testing}

Use grpcurl:

\begin{lstlisting}[language=bash]
# List services
grpcurl -plaintext localhost:9090 list

# Set velocity
grpcurl -plaintext -d '{"leftVelocity": 0.5, "rightVelocity": 0.5}' \
    localhost:9090 star.motor.MotorControl/SetVelocity

# Emergency stop
grpcurl -plaintext -d '{}' \
    localhost:9090 star.motor.MotorControl/EmergencyStop
\end{lstlisting}

% =============================================================================
\section{References}
% =============================================================================

\begin{itemize}
    \item Spring GraphQL: \url{https://docs.spring.io/spring-graphql/reference/}
    \item grpc-spring-boot-starter: \url{https://yidongnan.github.io/grpc-spring-boot-starter/}
    \item Pi4J v2: \url{https://pi4j.com/}
    \item Envoy gRPC-Web: \url{https://www.envoyproxy.io/docs/envoy/latest/configuration/http/http_filters/grpc_web_filter}
    \item STAR Software Architecture Primer: Section 10
\end{itemize}

